\documentclass[11pt]{article}
\usepackage{fontspec}
\usepackage[french,english]{babel}
\usepackage{amsmath}
\usepackage{amssymb}
\usepackage{amsthm}
\usepackage{mathrsfs}
\usepackage{tikz-cd}
\usepackage{enumitem}
\usepackage{lmodern}
\usetikzlibrary{babel}

\newcommand{\printedpage}[1]{\marginpar{\footnotesize p.~#1}}
\newcommand{\msp}[1]{\marginpar{\footnotesize [#1]}}
\newcommand{\uHom}{\underline{\mathrm{Hom}}}
\newcommand{\bDelta}{\mathbf{\Delta}}
\newcommand{\ho}{\mathrm{ho}}

\begin{document}

\printedpage{4}

\textbf{Cart (5) (Exactness axiom)} Let us consider two h-cartesian squares
\[
X = \left(\begin{tikzcd}
X_1 \ar[d] & X_3 \ar[l] \ar[d] \\
X_0 & X_2 \ar[l]
\end{tikzcd}\right)
\qquad \text{and} \qquad
Y = \left(\begin{tikzcd}
Y_1 \ar[d] & Y_3 \ar[l] \ar[d] \\
Y_0 & Y_2 \ar[l]
\end{tikzcd}\right)
\]
and a homomorphism $p = (p_0, p_1, p_2)$ between the diagrams of type
\[
\Phi \ = \
\begin{tikzcd}
1 \ar[d] & \\
0 & 2 \ar[l]
\end{tikzcd}
\]
associated with them (\textit{cf.} the solid-line diagram below):

\msp{page 5}
\[
\begin{tikzcd}[row sep=small, column sep=small]
& X_1 \ar[dl] \ar[dd, "p_1" near start] & & X_3 \ar[ll] \ar[dl] \ar[dd, dotted, "p_3"] \\
X_0 \ar[dd, "p_0"'] & & X_2 \ar[ll, crossing over] & \\
& Y_1 \ar[dl] & & Y_3 \ar[ll] \ar[dl] \\
Y_0 & & Y_2 \ar[ll] \ar[from=uu, crossing over, "p_2"' near end] & .
\end{tikzcd}
\]

Then:
\begin{itemize}[leftmargin=2.5em]
\item[a)] There exists a homomorphism $\bar{p} : X \longrightarrow Y$ inducing $p$, \textit{i.e.} there exists an arrow $p_3$ so as to complete the commutative cubical diagram (thus respecting the commutativity constraint for the two squares adjacent to the edge $p_3$).
\item[b)] If moreover $p_0$, $p_1$, $p_2$ are isomorphisms, the same is true of $p_3$.
\end{itemize}

Taking the case where $p_0$, $p_1$, $p_2$ are identities, we find:

\textbf{Corollary 1.} \textit{Under the conditions of} Cart (4)\textit{, if we have} two \textit{h-cartesian squares $X$, $X'$ which complete}
\[
\begin{tikzcd}[row sep=small]
& X_1 \ar[dl] & \\
X_0 & & , \\
& X_2 \ar[ul] &
\end{tikzcd}
\]
\textit{there exists a homomorphism from $X$ to $X'$ which induces the identity on the vertices with indices} 0, 1, 2\textit{, and such a homomorphism is always an isomorphism.}

But this isomorphism is not necessarily unique, in other words, an automorphism of an h-cartesian square can induce the identity on the vertices $X_0$, $X_1$, $X_2$, without being the identity.

\printedpage{5}

\msp{page 6}
\textbf{Corollary 2.} \textit{Let us consider a commutative square}
\[
Y = \left(\begin{tikzcd}
Y_1 \ar[d] & Y_3 \ar[l] \ar[d] \\
Y_0 & Y_2 \ar[l]
\end{tikzcd}\right) \quad ,
\]
\textit{and an object $T$ of $\mathcal{H}$, whence (putting} $\mathrm{Hom}(T,X) = X(T)$\textit{) a commutative square in} (Ens)
\[
\begin{tikzcd}
Y_1(T) \ar[d] & Y_3(T) \ar[l] \ar[d] \\
Y_0(T) & Y_2(T) \ar[l]
\end{tikzcd} \quad ,
\]
\textit{and hence}
\[
Y_3(T) \longrightarrow Y_1(T) \times_{Y_0(T)} Y_2(T) \quad . \tag*{$(*)$}
\]
\textit{If $Y$ is h-cartesian, this map} $(*)$ \textit{is} surjective \textit{whatever $T$ may be.}

We apply Cart (5) to $Y$ and to the constant square $X$ with value $T$
\[
X = \quad
\begin{tikzcd}
T \ar[d, "\mathrm{id}"'] & T \ar[l, "\mathrm{id}"'] \ar[d, "\mathrm{id}"] \\
T & T \ar[l, "\mathrm{id}"']
\end{tikzcd} \quad ,
\]
noting that giving $p = (p_0, p_1, p_2)$ from
\[
\begin{tikzcd}
T \ar[d] & \\
T & T \ar[l]
\end{tikzcd}
\qquad\qquad \text{to} \qquad\qquad
\begin{tikzcd}
Y_1 \ar[d] & \\
Y_0 & Y_2 \ar[l]
\end{tikzcd}
\]
amounts to giving an element of the right-hand side of $(*)$,
\[
\begin{tikzcd}
Y_1 \ar[d] & T \ar[l, "p_1"'] \ar[dl, "p_0"] \ar[d, "p_2"] \\
Y_0 & Y_2 \ar[l]
\end{tikzcd} \quad ,
\]
and that the elements of $Y_3(T)$ which map (by $(*)$) to this element are none other than the $p_3$ whose existence axiom Cart (5) asserts.

\textbf{N.B.} We shall note that the surjective map $(*)$ is in general \textit{not} bijective, \textit{i.e.} the h-cartesian square is in general not cartesian.

\printedpage{6}

\msp{page 7}
\textbf{N.B.} There is a dual notion of \textit{category with h-cocartesian squares,} or \textit{right quadrangulated category.} Most often, the two structures are present simultaneously. We should bring out properties which relate the two dual structures on one and the same category. Thus, unless I am mistaken, in the case $\mathcal{H}$ additive, the two kinds of squares, the h-cartesian and the h-cocartesian ones, should \textit{coincide.} (I shall have to come back to this later.) In the case $\mathcal{H}$ ``pointed'', there is the adjunction of the functors $\Sigma$, $\Omega$ linked, one to the right quadrangulated structure, the other to the left quadrangulated structure. And there are, when there are internal or external \underline{Hom}, possible properties of $\uHom(X,Y)$ with respect to h-cocartesian quadrangles in $X$, h-cartesian in $Y$ -- but this is hardly a link between the two structures \ldots

\section*{2\quad Examples}

A first trivial example, is the one where $\mathcal{H}$ is stable under fiber products, and where the h-cartesian squares are defined as the cartesian squares. Axioms Cart (1) to Cart (5) are then well-known, $\pm$ tautological properties of the notion of a cartesian square. And
\msp{page 8}
in a sense, the general case must reduce to this banal special case.

The only type of other case that I know of, is the one where $\mathcal{H} = \mathcal{M}W^{-1}$, where $W$ is a localizer admitting a theory of the $W$-homotopy fiber (\textit{cf.} chapter XII). In this case, as we have said, we take as h-cartesian squares in $\mathcal{H}$ those which are isomorphic, in $\uHom(\bDelta_1 \times \bDelta_1, \mathcal{H})$, to a $W$-cartesian square in $\mathcal{M}$. We still have to verify axioms Cart (1) to Cart (5). I think that I have little chance of getting there except by means of a

\textbf{Proposition 1.} \textit{Under hypotheses on $(\mathcal{M}, W)$ which will be brought out below, satisfied in any case if $(\mathcal{M}, W)$ is Quillenizable, on the category of finite ordered sets} (\textit{cf.} XIII, p. 51)\textit{, we have this: let}
\[
X = \left(\begin{tikzcd}
X_1 \ar[d] & X_3 \ar[l] \ar[d] \\
X_0 & X_2 \ar[l]
\end{tikzcd}\right)
\]
\textit{be a square in $\mathcal{M}$, and let} $\ho_W(X)$ \textit{be its image in} $\uHom(\bDelta_1 \times \bDelta_1, \mathcal{H})$\textit{. For} $\ho_W(X)$ \textit{to be h-cartesian, it is} necessary (\textit{and it is sufficient, of course}) \textit{that $X$ be $W$-cartesian} \textit{.}

\printedpage{7}

% the editors' afterword (Postface des éditeurs) is not reproduced (their own writing)

\end{document}
